% !TeX root = closed-systems-from-comparison-completeness.tex
% ============================================================
\chapter{Conditional Numerical Closure: Quartic Carrier and Mass--Depth Law}
\label{ch:matter-numbers}
% ============================================================

% ============================================================
\begin{chapterorientation}
\orientationitem{Input}{The scalar invariant \(m_0(\sigma)\), refinement depth \(k_\sigma\), and sector-separation results of \cref{ch:matter}.}
\orientationitem{Role}{This chapter develops the conditional numerical closure layer carried by quartic depth structure.}
\orientationitem{Output}{The quartic mass-coordinate law, anchor and ratio geometry, and an integrated conditional closure theorem.}
\end{chapterorientation}
\section{Introduction and logical status}
\label{sec:numbers-intro}
% ============================================================

Chapter~\ref{ch:matter} proves that primitive sectors are classified by
internal gauge representation data together with depth information
\(k_\sigma\) (\cref{def:spectral-depth,thm:spectral-depth-separation}),
and that the quadratic scalar invariant
\(m_0(\sigma)=Q(\kappa_\sigma)\) is transport-invariant and strictly
positive (\cref{thm:spectral-scalar-positive,thm:spectral-scalar-invariant}).
It also records that \(m_0\) is not injective on sectors
(\cref{rem:m0-not-injective}).

This chapter develops a quartic refinement of the quadratic carrier:
an invariant depth observable constructed from degree--\(4\) combinatorics.
The guiding objective is to separate theorem-level identities from
external calibration data, so that each numerical statement is
explicitly typed as either a structural consequence or boundary data.

The argument is organized in four components.
\begin{enumerate}[label=\textup{(\roman*)}]
\item A rigidity envelope for admissible scalar depth observables is proved:
      under intrinsicality and bilinear-origin hypotheses, the only universal
      leading law is quartic.
\item A canonical depth profile \(\eta_k\in\Lambda^2(\mathbb Z^k)\) is
      introduced, and its quartic square is computed exactly.
\item A quartic depth observable \(\mathfrak m_4\) is defined and shown to
      satisfy the exact law
      \(\mathfrak m_4(\sigma)=4\mu\binom{k_\sigma}{4}\).
\item External data required for anchor/depth calibration (explicit or
      recovered), together with cubic selection,
      are stated explicitly.
\end{enumerate}

Because these numerical identifications are not fixed by the structural
closure chain, they are treated as explicit \emph{boundary data} in the
sense of \cref{app:conditional-completion,sec:cond-open}.
No theorem-level claim is used beyond proved structure plus those stated
data.

% ============================================================
\section{Scalar rigidity envelope for depth observables}
\label{sec:numbers-rigidity-envelope}
% ============================================================

The attached notes package contributes a useful meta-level rigidity claim:
before fixing any explicit formula, admissible scalar depth observables
can already be constrained by intrinsicality and transport covariance.
This section records that envelope in theorem form and then specializes
it to the quartic law used in the remainder of the chapter.

\begin{definition}[Admissible scalar depth observable]
\label{def:numbers-admissible-depth-observable}
A map \(\Phi:\mathbb N\to\mathbb R\) is called an
\emph{admissible scalar depth observable} if:
\begin{enumerate}[label=\textup{(\roman*)}]
\item
      \textbf{Intrinsicality:} \(\Phi(k)\) depends only on the depth-\(k\)
      quadratic carrier and canonical scalar channel, with no auxiliary
      coordinate or embedding choice;
\item
      \textbf{Transport invariance:} \(\Phi\) is invariant under admissible
      transport/gauge identifications on the carrier;
\item
      \textbf{Quadratic-bilinear origin:} \(\Phi(k)\) is obtained from a
      scalar contraction of a bilinear expression on the quadratic carrier;
\item
      \textbf{Uniform depth covariance:} the assignment
      \(k\mapsto\Phi(k)\) is functorial under canonical depth inclusions and
      polynomially bounded in \(k\).
\end{enumerate}
\end{definition}

\begin{proposition}[Quartic growth envelope]
\label{prop:numbers-quartic-envelope}
For every admissible scalar depth observable
\(\Phi\) in the sense of
\cref{def:numbers-admissible-depth-observable}, there exists a
constant \(C>0\) such that
\[
  |\Phi(k)|\le C(1+k^4)
  \qquad (k\in\mathbb N).
\]
\end{proposition}

\begin{proof}
By \cref{thm:spectral-binomial}, the quadratic carrier rank is
\(\dim\mathcal K_k=\binom{k}{2}=O(k^2)\). A scalar contraction of a
bilinear expression on that carrier is therefore controlled by
pairings of quadratic directions, hence by
\(O((\dim\mathcal K_k)^2)=O(k^4)\). This is the claimed quartic
growth envelope.
\end{proof}

\begin{theorem}[Quartic leading-order rigidity under scalar descent]
\label{thm:numbers-pure-quartic-rigidity}
Assume \(\Phi\) satisfies
\cref{def:numbers-admissible-depth-observable} and additionally:
\begin{enumerate}[label=\textup{(\roman*)}]
\item \(\Phi(0)=0\);
\item odd contributions are eliminated by orientation/phase descent;
\item the leading asymptotic growth is extensive in
      \(\binom{k}{2}^2\).
\end{enumerate}
Then there exists \(\alpha>0\) such that
\[
  \Phi(k)=\alpha k^4+o(k^4)
  \qquad (k\to\infty).
\]
In particular, the universal leading law is quartic.
\end{theorem}

\begin{proof}
By \cref{prop:numbers-quartic-envelope}, \(\Phi(k)=O(k^4)\). The
hypothesized extensivity in \(\binom{k}{2}^2\) identifies the leading
asymptotic behavior with a positive multiple of
\(\binom{k}{2}^2\sim \tfrac14 k^4\), so there exists \(\alpha>0\)
such that \(\Phi(k)=\alpha k^4+o(k^4)\). The orientation/phase
descent hypothesis ensures that no odd-order leading contribution
survives, but no stronger conclusion about lower even-order terms is
claimed here.
\end{proof}

\begin{remark}[Role of the cubic selector]
\label{rem:numbers-rigidity-vs-selector}
The rigidity theorem isolates the universal leading behavior of the
depth law.
It does \emph{not} derive a selector for admissible depths.
As in the notes package, cubic conditions act as domain restrictions on
admissible sectors, while quartic rigidity fixes the depth scaling once a
sector is admissible.
\end{remark}

% ============================================================
\section{Depth profiles in the quadratic carrier}
\label{sec:numbers-depth-profile}
% ============================================================

By \cref{thm:spectral-binomial,thm:binomk2}, the stage-\(k\) quadratic
carrier is identified with
\(\Lambda^2(\mathbb Z^k)\).
Fix the ordered basis \(e_1,\dots,e_k\) of \(\mathbb Z^k\).

\begin{definition}[Cumulative depth profile]
\label{def:numbers-depth-profile}
For \(k\in\mathbb N\), define
\[
  \eta_k
  :=
  \sum_{1\le i<j\le k} e_i\wedge e_j
  \in \Lambda^2(\mathbb Z^k).
\]
\end{definition}

\begin{remark}[Quadratic count compatibility]
\label{rem:numbers-q2-compatibility}
The number of summands in \(\eta_k\) is \(\binom{k}{2}\),
matching the quadratic channel count of \cref{thm:binomk2} and the
normalization variable used in \cref{thm:forced-normalization}.
\end{remark}

\begin{definition}[Quartic profile]
\label{def:numbers-quartic-profile}
Define
\[
  \omega_k := \eta_k\wedge\eta_k
  \in \Lambda^4(\mathbb Z^k).
\]
\end{definition}

\begin{lemma}[Explicit quartic expansion]
\label{lem:numbers-quartic-expansion}
For every \(k\ge 4\),
\[
  \omega_k
  =
  2\sum_{1\le i_1<i_2<i_3<i_4\le k}
  e_{i_1}\wedge e_{i_2}\wedge e_{i_3}\wedge e_{i_4}.
\]
\end{lemma}

\begin{proof}
Write \(u_{ab}:=e_a\wedge e_b\) for \(a<b\).
Then
\[
  \omega_k
  =
  \sum_{a<b}\sum_{c<d} u_{ab}\wedge u_{cd}.
\]
If \(\{a,b\}\cap\{c,d\}\neq\varnothing\), the wedge term vanishes.
Hence only disjoint pairs contribute.

Fix \(i_1<i_2<i_3<i_4\).
The nonzero contributions using exactly these four indices come from the
three unordered pairings:
\[
  (i_1i_2)(i_3i_4),\qquad
  (i_1i_3)(i_2i_4),\qquad
  (i_1i_4)(i_2i_3).
\]
Relative to the oriented basis element
\(e_{i_1}\wedge e_{i_2}\wedge e_{i_3}\wedge e_{i_4}\), their signs are
\(+1,-1,+1\), so the signed sum is \(1\).
Because degree--\(2\) forms commute under wedge
(\(\alpha\wedge\beta=\beta\wedge\alpha\) for \(\deg\alpha=\deg\beta=2\)),
each unordered pairing appears twice in the ordered double sum.
Therefore the total coefficient of each ordered \(4\)-tuple is
\(2\cdot 1=2\), proving the formula.
\end{proof}

\begin{definition}[Quartic coefficient energy]
\label{def:numbers-quartic-energy}
For
\(\omega=\sum_I c_I e_I\in\Lambda^4(\mathbb R^k)\), where
\(I=(i_1<i_2<i_3<i_4)\) and
\(e_I:=e_{i_1}\wedge e_{i_2}\wedge e_{i_3}\wedge e_{i_4}\), define
\[
  E_4(\omega):=\sum_I c_I^2.
\]
\end{definition}

\begin{proposition}[Quartic counting law]
\label{prop:numbers-quartic-count}
For \(k\ge 4\),
\[
  E_4(\omega_k)=4\binom{k}{4}.
\]
\end{proposition}

\begin{proof}
By \cref{lem:numbers-quartic-expansion}, every basis \(4\)-blade appears
with coefficient \(2\), and there are exactly \(\binom{k}{4}\) such blades.
Hence
\[
  E_4(\omega_k)
  =
  \binom{k}{4}\cdot 2^2
  =
  4\binom{k}{4}.
\]
\end{proof}

% ============================================================
\section{Quartic depth observable and mass-coordinate law}
\label{sec:numbers-mass-law}
% ============================================================

\begin{definition}[Quartic depth observable]
\label{def:numbers-m4}
Fix \(\mu>0\).
For a nontrivial primitive sector \(\sigma\) with depth \(k_\sigma\)
(\cref{def:spectral-depth}), define
\[
  \mathfrak m_4(\sigma)
  :=
  \mu\,E_4\!\bigl(\omega_{k_\sigma}\bigr)
  =
  \mu\,E_4\!\bigl(\eta_{k_\sigma}\wedge\eta_{k_\sigma}\bigr).
\]
\end{definition}

\begin{theorem}[Exact quartic formula]
\label{thm:numbers-quartic-formula}
For every nontrivial primitive sector \(\sigma\),
\[
  \mathfrak m_4(\sigma)=4\mu\binom{k_\sigma}{4}.
\]
\end{theorem}

\begin{proof}
Combine \cref{def:numbers-m4} with
\cref{prop:numbers-quartic-count} evaluated at \(k=k_\sigma\).
\end{proof}

\begin{theorem}[Intrinsic quartic mass-coordinate bridge]
\label{thm:numbers-mass-calibration}
For every nontrivial primitive sector \(\sigma\), the closed-system
mass coordinate determined by the quartic coefficient-energy readout is
\[
  m_\sigma:=\mathfrak m_4(\sigma).
\]
\end{theorem}

\begin{proof}
The closed-system scalar readout available at quartic depth is the
coefficient energy of the canonical quartic profile
\(\omega_{k_\sigma}\), normalized by the universal positive scale
\(\mu\).  By \cref{def:numbers-m4}, that readout is exactly
\[
  \mu E_4(\omega_{k_\sigma})=\mathfrak m_4(\sigma).
\]
Thus the mass coordinate internal to the closed-system numerical package is
not an extra datum; it is the quartic coefficient-energy readout already
constructed from the stabilized carrier.
\end{proof}

\begin{corollary}[Mass-coordinate law]
\label{cor:numbers-conditional-mass-law}
For every nontrivial primitive sector,
\[
  m_\sigma=4\mu\binom{k_\sigma}{4}.
\]
\end{corollary}

\begin{proof}
Immediate from \cref{thm:numbers-quartic-formula}.
\end{proof}

\begin{remark}[What is forced vs. what is calibrated]
\label{rem:numbers-forced-vs-calibrated}
The combinatorial identity
\(\mathfrak m_4(\sigma)=4\mu\binom{k_\sigma}{4}\)
is theorem-level once \(\mathfrak m_4\) is defined.
The closed-system mass coordinate is therefore theorem-level.  Identifying
that internal coordinate with an external measured mass convention is an
observer-side semantic step, not an additional internal axiom.
\end{remark}

% ============================================================
\section{Depth anchors and ratio package}
\label{sec:numbers-ratios}
% ============================================================

\begin{definition}[Explicit depth-anchor package]
\label{def:numbers-depth-anchors}
The explicit \(6,18,35\) depth-anchor package is the declaration that the
three reference sectors \(e,\mu,\tau\) have depths
\[
  k_e=6,
  \qquad
  k_\mu=18,
  \qquad
  k_\tau=35.
\]
\end{definition}

\begin{corollary}[Conditional ratio predictions]
\label{cor:numbers-ratios}
For the closed-system mass coordinate and the explicit anchor package
of \cref{def:numbers-depth-anchors},
\[
  \frac{m_\mu}{m_e}
  =
  \frac{\binom{18}{4}}{\binom{6}{4}}
  =204,
\]
\[
  \frac{m_\tau}{m_e}
  =
  \frac{\binom{35}{4}}{\binom{6}{4}}
  =\frac{10472}{3}
  \approx 3490.67.
\]
\end{corollary}

\begin{proof}
In mass ratios, the normalization constant \(4\mu\) cancels.
Substitute the depth values from
\cref{def:numbers-depth-anchors} into
\cref{cor:numbers-conditional-mass-law}.
\end{proof}

% ============================================================
\section{Cubic obstruction from phase typing (optional selector data)}
\label{sec:numbers-cubic}
% ============================================================

This section introduces an internal phase-typing selector for admissible
depth profiles.
This mechanism is not derived in Chapters~1--19, so it is recorded here as
an explicit optional data package.

\begin{definition}[Phase data at depth \(k\)]
\label{def:numbers-phase-data}
Fix a charge unit \(q_0\in\mathbb N_{>0}\).
For a depth-\(k\) profile choose integer labels
\(q_1,\dots,q_k\in\mathbb Z\), and define phase factors
\[
  \phi_i
  :=
  \exp\!\left(\frac{2\pi i}{q_0}\,q_i\right)
  \qquad (1\le i\le k).
\]
\end{definition}

\begin{definition}[Phase admissibility]
\label{def:numbers-phase-admissible}
A depth-\(k\) phase configuration is \emph{admissible} if it extends to a
coherent phase assignment on all higher refinement levels.
\end{definition}

\begin{lemma}[Triple coherence]
\label{lem:numbers-triple-coherence}
In this optional phase-typing model, admissibility requires triviality of
the \(3\)-cocycle
\[
  \Phi(i,j,\ell)
  :=
  \exp\!\left(\frac{2\pi i}{q_0^3}\,q_iq_jq_\ell\right).
\]
Consequently,
\[
  \sum_{1\le i<j<\ell\le k} q_iq_jq_\ell=0.
\]
\end{lemma}

\begin{proof}
Refinement invariance forces associativity of phase composition on triple
overlaps.
This produces a \(3\)-cocycle obstruction represented by
\(\Phi(i,j,\ell)\), and admissibility requires that obstruction class to be
trivial.
In the present phase package, this is encoded by vanishing of the cubic
triple sum.
\end{proof}

\begin{definition}[Cubic charge defect]
\label{def:numbers-cubic-defect}
Given integer labels \(q_1,\dots,q_k\), define
\[
  \Delta_3(k):=\sum_{i=1}^k q_i^3.
\]
\end{definition}

\begin{lemma}[Cubic obstruction under charge neutrality]
\label{lem:numbers-cubic-obstruction}
Assume
\[
  \sum_{i=1}^k q_i=0.
\]
If
\[
  \sum_{1\le i<j<\ell\le k} q_iq_jq_\ell=0,
\]
then
\[
  \Delta_3(k)=0.
\]
\end{lemma}

\begin{proof}
Let \(p_r:=\sum_i q_i^r\) and \(e_r\) be the elementary symmetric sums.
Newton's identity gives
\[
  p_3=e_1^3-3e_1e_2+3e_3.
\]
Under charge neutrality, \(e_1=p_1=\sum_i q_i=0\), so \(p_3=3e_3\).
Therefore \(e_3=\sum_{i<j<\ell}q_iq_jq_\ell=0\) implies
\(p_3=\sum_i q_i^3=0\), which is exactly \(\Delta_3(k)=0\).
\end{proof}

\begin{corollary}[Cubic selector induced by phase typing]
\label{cor:numbers-cubic-criterion}
Under phase admissibility and charge neutrality,
\[
  \Delta_3(k)=0
\]
is a necessary condition for depth-\(k\) admissibility.
\end{corollary}

\begin{proof}
Combine \cref{lem:numbers-triple-coherence,lem:numbers-cubic-obstruction}.
\end{proof}

\begin{proposition}[Minimal depth from cubic cancellation]
\label{prop:numbers-cubic-k35}
Assume \cref{cor:numbers-cubic-criterion} and suppose
\[
  \Delta_3(34)\neq0,
  \qquad
  \Delta_3(35)=0.
\]
Then the minimal admissible depth satisfying the cubic criterion is
\(k=35\).
\end{proposition}

\begin{proof}
By \cref{cor:numbers-cubic-criterion}, admissibility at depth \(k\)
requires \(\Delta_3(k)=0\).
The given values rule out \(k=34\) and admit \(k=35\).
Hence the least admissible depth under this criterion is \(35\).
\end{proof}

\begin{corollary}[Conditional third-generation depth]
\label{cor:numbers-k-tau-35}
If the third-generation onset depth is identified with the first depth
admissible under the cubic phase-typing criterion, then \(k_\tau=35\).
\end{corollary}

% ============================================================
\section{Structural dependency map with the closure chain}
\label{sec:numbers-dependency-map}
% ============================================================

For traceability against the preceding chapters, this section records the
exact structural interfaces used by the numerical analysis.

\begin{definition}[Dependency interface]
\label{def:numbers-dependency-interface}
The Chapter~20 dependency interface consists of the following four dependencies:
\begin{enumerate}[label=\textup{(\roman*)}]
\item
      \textbf{Depth typing and invariance.}
      Primitive sectors carry a depth index
      \(k_\sigma\) that is transport-invariant
      (\cref{def:spectral-depth,prop:spectral-depth-invariant}).
\item
      \textbf{Depth separation and hierarchy.}
      Distinct depth values represent distinct hierarchy levels
      (\cref{thm:spectral-depth-separation,cor:spectral-hierarchy}).
\item
      \textbf{Quadratic carrier structure.}
      The stabilized degree--\(2\) carrier is identified with
      \(\Lambda^2(\mathbb Z^k)\), with rank \(\binom{k}{2}\)
      (\cref{thm:spectral-binomial,thm:binomk2}).
\item
      \textbf{Normalization ledger.}
      Quantitatively open entries are tracked as boundary data
      (\cref{app:conditional-completion,sec:cond-open}).
\end{enumerate}
\end{definition}

\begin{proposition}[Logical role of the dependency interface]
\label{prop:numbers-dependency-role}
The unconditional quartic-combinatorics package in this chapter is a
formal consequence of
\cref{def:numbers-dependency-interface} and multilinear algebra on
\(\Lambda^2(\mathbb Z^k)\).
The scalar rigidity-envelope results and the optional cubic-selector
results use additional chapter-local hypotheses, recorded explicitly in
\cref{def:numbers-admissible-depth-observable,def:numbers-phase-admissible}.
Concrete generation-depth calibration is conditional on explicit calibration
data: either fixed anchor depths
\cref{def:numbers-depth-anchors}, or the ratio/cubic recovery route of
\cref{thm:numbers-anchor-closure-61835}.
\end{proposition}

\begin{proof}
Depth labels \(k_\sigma\) and their invariance enter through
\cref{def:spectral-depth,prop:spectral-depth-invariant}, and the
carrier used for quartic combinatorics is fixed by
\cref{thm:spectral-binomial,thm:binomk2}.
Therefore the identities involving \(\eta_k\), \(\omega_k\),
coefficient counts, the polynomial \(q(k)=4\binom{k}{4}\), and their
derived ratio, monotonicity, inversion, and scaling consequences are
structural.
The scalar rigidity-envelope package additionally uses the local
admissibility hypotheses of
\cref{def:numbers-admissible-depth-observable}, and the optional cubic
selector package uses the separately stated phase-admissibility and
charge-neutrality hypotheses recorded in
\cref{def:numbers-phase-admissible,lem:numbers-cubic-obstruction,thm:numbers-cubic-equivalent-forms}.
The transition from these structural or locally hypothesized identities
to closed-system mass coordinates uses the theorem-level bridge
\(m_\sigma=\mathfrak m_4(\sigma)\)
(\cref{thm:numbers-mass-calibration}), and concrete anchor
calibration uses either explicit depth anchors
(\cref{def:numbers-depth-anchors}) or the ratio/cubic recovery route
(\cref{thm:numbers-anchor-closure-61835}).
The observer-side interpretation of these closed-system coordinates is
tracked separately in \cref{sec:cond-open}.
Hence the claimed logical split follows.
\end{proof}

\begin{definition}[Chapter~20 numerical output types]
\label{def:numbers-output-types}
A Chapter~20 numerical output is any statement of one of the following
forms:
\begin{enumerate}[label=\textup{(O\arabic*)}]
\item an exact quartic coefficient-count or polynomial identity for
      \(E_4(\omega_k)\), \(q(k)=4\binom{k}{4}\), or their finite
      differences;
\item a closed-system mass-coordinate identity obtained from
      \(m_\sigma=\mathfrak m_4(\sigma)\);
\item an order, convexity, scaling, or finite-inversion consequence of
      \(q(k)\);
\item an anchor-ratio consequence for a specified depth-anchor package or
      for the ratio/cubic recovery route;
\item a cubic-selector consequence using the separately specified
      phase-typing data.
\end{enumerate}
\end{definition}

\begin{theorem}[Numerical dependency normal form]
\label{thm:numbers-dependency-normal-form}
Every Chapter~20 numerical output factors through the ordered dependency
tuple
\[
  \mathcal D_{20}
  :=
  \bigl(
    k_\sigma,
    \Lambda^2(\mathbb Z^{k_\sigma}),
    \eta_{k_\sigma},
    \mu,
    \mathcal A,
    \Delta_3
  \bigr),
\]
where \(\mathcal A\) denotes either an explicitly specified anchor package
or the ratio/cubic recovery data, and \(\Delta_3\) denotes optional
cubic-selector data.  More precisely:
\begin{enumerate}[label=\textup{(\roman*)}]
\item outputs of type \textup{(O1)} use only
      \((k_\sigma,\Lambda^2(\mathbb Z^{k_\sigma}),\eta_{k_\sigma})\);
\item outputs of types \textup{(O2)} and \textup{(O3)} use those structural
      data together with the positive scale \(\mu\);
\item outputs of type \textup{(O4)} additionally use exactly the anchor data
      named in the relevant statement;
\item outputs of type \textup{(O5)} additionally use exactly the
      phase-typing data named in the relevant statement.
\end{enumerate}
No further numerical datum is used implicitly.
\end{theorem}

\begin{proof}
For \textup{(O1)}, the coefficient computation is
\cref{prop:numbers-quartic-count,cor:numbers-energy-by-counting}; its proof
uses only the canonical exterior-square carrier and the profile \(\eta_k\).
The difference tower, recurrence, monotonicity, convexity, and elementary
binomial estimates are then algebraic consequences of the same polynomial
\(q(k)=4\binom{k}{4}\).

For \textup{(O2)}, \cref{thm:numbers-mass-calibration} identifies the
closed-system mass coordinate with \(\mathfrak m_4\), and
\cref{cor:numbers-conditional-mass-law} substitutes the exact quartic formula;
the only additional scalar datum is the positive normalization \(\mu\) already
present in \cref{def:numbers-m4}.  Type \textup{(O3)} statements are obtained
from the same formula by applying the order and inversion properties of
\(q(k)\), so no new datum enters.

For \textup{(O4)}, \cref{thm:numbers-general-ratio-law} cancels the common
factor \(4\mu\) and leaves only the specified anchor depths.  When those
depths are not stipulated, \cref{thm:numbers-anchor-closure-61835} supplies
them from the explicitly named ratio/cubic recovery data.  For
\textup{(O5)}, the cubic statements cite
\cref{def:numbers-phase-admissible,cor:numbers-cubic-criterion} and therefore
use exactly the phase-typing data displayed there.

These cases exhaust \cref{def:numbers-output-types}.  Each proof in the
chapter cites one of the structural, mass-coordinate, anchor, or cubic routes
listed above; hence any dependency not appearing in \(\mathcal D_{20}\) would
have no cited entry point.  Therefore no further numerical datum is used
implicitly.
\end{proof}

\begin{corollary}[Auditability of Chapter~20 outputs]
\label{cor:numbers-auditability}
To audit any Chapter~20 numerical claim, it is enough to identify its output
type in \cref{def:numbers-output-types} and then read off the corresponding
component of \(\mathcal D_{20}\) from
\cref{thm:numbers-dependency-normal-form}.
\end{corollary}

\begin{proof}
The theorem gives an exhaustive case split by output type and records the
exact dependency components used in each case.  Thus checking the output type
determines the complete dependency list.
\end{proof}

\begin{remark}[Consistency with Chapter~19 scope]
\label{rem:numbers-scope-ch19}
Chapter~19 establishes positivity and transport invariance for
\(m_0\), together with non-injectivity of \(m_0\)
(\cref{thm:spectral-scalar-positive,thm:spectral-scalar-invariant,rem:m0-not-injective}).
The quartic observable \(\mathfrak m_4\) introduced here is not asserted
to replace these structural facts;
it is an additional depth-based statistic with separate conditional
interpretation.
\end{remark}

% ============================================================
\section{Coefficient-level decomposition of the quartic profile}
\label{sec:numbers-coeff-decomposition}
% ============================================================

This section expands \cref{lem:numbers-quartic-expansion} into a fully
indexed coefficient computation.

\begin{definition}[Support of a wedge monomial]
\label{def:numbers-support}
For \(a<b\), define
\(
  u_{ab}:=e_a\wedge e_b
\).
For \((a,b,c,d)\) with \(a<b\), \(c<d\), define the support set
\[
  \supp(a,b;c,d):=\{a,b,c,d\}.
\]
\end{definition}

\begin{lemma}[Support criterion for nonzero terms]
\label{lem:numbers-support-criterion}
For \(a<b\), \(c<d\), the wedge term
\(u_{ab}\wedge u_{cd}\) is nonzero if and only if
\(|\supp(a,b;c,d)|=4\).
\end{lemma}

\begin{proof}
If \(\supp(a,b;c,d)\) has cardinality strictly smaller than \(4\), then one
index appears at least twice in
\(e_a\wedge e_b\wedge e_c\wedge e_d\), hence alternatingness of exterior
product implies
\(u_{ab}\wedge u_{cd}=0\).

Conversely, if \(|\supp(a,b;c,d)|=4\), then
\(e_a,e_b,e_c,e_d\) are distinct basis vectors, so their wedge is a nonzero
multiple of the oriented basis element determined by the sorted index order.
Thus \(u_{ab}\wedge u_{cd}\neq0\).
\end{proof}

\begin{definition}[Coefficient extraction]
\label{def:numbers-coefficient-extraction}
For each ordered \(4\)-subset
\(I=(i_1<i_2<i_3<i_4)\), define
\(e_I:=e_{i_1}\wedge e_{i_2}\wedge e_{i_3}\wedge e_{i_4}\).
If
\(\omega=\sum_I c_I e_I\), denote by
\(\operatorname{coeff}_I(\omega):=c_I\) the \(I\)-coefficient.
\end{definition}

\begin{lemma}[Pairing-sign table on a fixed support]
\label{lem:numbers-pairing-sign-table}
Fix \(I=(i_1<i_2<i_3<i_4)\).
Relative to \(e_I\), the three unordered pairings satisfy
\[
  u_{i_1i_2}\wedge u_{i_3i_4}=+e_I,
\]
\[
  u_{i_1i_3}\wedge u_{i_2i_4}=-e_I,
\]
\[
  u_{i_1i_4}\wedge u_{i_2i_3}=+e_I.
\]
Hence their signed sum equals \(e_I\).
\end{lemma}

\begin{proof}
The first identity is immediate from ordering.
For the second,
\[
  e_{i_1}\wedge e_{i_3}\wedge e_{i_2}\wedge e_{i_4}
  =-e_{i_1}\wedge e_{i_2}\wedge e_{i_3}\wedge e_{i_4}
  =-e_I,
\]
since one transposition exchanges \(i_3\) and \(i_2\).
For the third,
\[
  e_{i_1}\wedge e_{i_4}\wedge e_{i_2}\wedge e_{i_3}
  =(+1)e_I,
\]
because the permutation
\((1,4,2,3)\to(1,2,3,4)\) has two inversions.
Summing gives the claim.
\end{proof}

\begin{proposition}[Uniform coefficient law]
\label{prop:numbers-uniform-coefficient}
For every ordered \(4\)-subset \(I\), one has
\[
  \operatorname{coeff}_I(\omega_k)=2.
\]
Equivalently,
\[
  \omega_k=2\sum_I e_I,
\]
where \(I\) runs through all ordered \(4\)-subsets of
\(\{1,\dots,k\}\).
\end{proposition}

\begin{proof}
Expand
\(
  \omega_k=\sum_{a<b}\sum_{c<d}u_{ab}\wedge u_{cd}
\).
By \cref{lem:numbers-support-criterion}, only terms with four distinct
indices contribute.
Fix \(I\).
Within support \(I\), the unordered contribution equals \(e_I\) by
\cref{lem:numbers-pairing-sign-table}.
Because degree--\(2\) forms commute,
\(
  u_{ab}\wedge u_{cd}=u_{cd}\wedge u_{ab}
\),
each unordered pairing appears twice in the ordered double sum.
Hence \(\operatorname{coeff}_I(\omega_k)=2\).
Summing over \(I\) yields the displayed expansion.
\end{proof}

\begin{corollary}[Energy identity by coefficient counting]
\label{cor:numbers-energy-by-counting}
For \(k\ge4\),
\[
  E_4(\omega_k)
  =
  \sum_I \operatorname{coeff}_I(\omega_k)^2
  =
  \sum_I 4
  =
  4\binom{k}{4}.
\]
\end{corollary}

\begin{proof}
Apply \cref{prop:numbers-uniform-coefficient} and count the number of
ordered \(4\)-subsets.
\end{proof}

% ============================================================
\section{Polynomial form and discrete quartic calculus}
\label{sec:numbers-discrete-calculus}
% ============================================================

\begin{definition}[Reduced quartic depth polynomial]
\label{def:numbers-qk}
Define
\[
  q(k):=4\binom{k}{4},
\]
so that
\(
  \mathfrak m_4(\sigma)=\mu\,q(k_\sigma)
\)
by \cref{thm:numbers-quartic-formula}.
\end{definition}

\begin{proposition}[Closed polynomial form]
\label{prop:numbers-qk-polynomial}
For all integers \(k\),
\[
  q(k)=\frac{1}{6}k(k-1)(k-2)(k-3).
\]
\end{proposition}

\begin{proof}
By definition,
\[
  q(k)=4\binom{k}{4}=4\frac{k(k-1)(k-2)(k-3)}{24},
\]
which simplifies to the displayed expression.
\end{proof}

\begin{lemma}[Vanishing order at low depth]
\label{lem:numbers-qk-vanishing}
One has
\[
  q(0)=q(1)=q(2)=q(3)=0,
  \qquad
  q(4)=4.
\]
\end{lemma}

\begin{proof}
The factorized form from \cref{prop:numbers-qk-polynomial} contains the
four consecutive factors
\(k\), \(k-1\), \(k-2\), \(k-3\), so the first four values vanish.
At \(k=4\),
\(
  q(4)=4\binom{4}{4}=4
\).
\end{proof}

\begin{definition}[Forward difference operator]
\label{def:numbers-forward-diff}
For a function \(f:\mathbb Z\to\mathbb R\), define
\[
  (\Delta f)(k):=f(k+1)-f(k).
\]
Inductively,
\(
  \Delta^{r+1}f:=\Delta(\Delta^r f)
\).
\end{definition}

\begin{theorem}[Finite-difference tower for \(q\)]
\label{thm:numbers-difference-tower}
For every integer \(k\), one has
\[
  \Delta q(k)=4\binom{k}{3},
\]
\[
  \Delta^2 q(k)=4\binom{k}{2},
\]
\[
  \Delta^3 q(k)=4k,
\]
\[
  \Delta^4 q(k)=4,
\]
\[
  \Delta^5 q(k)=0.
\]
\end{theorem}

\begin{proof}
Use Pascal's identity
\(
  \binom{k+1}{r}-\binom{k}{r}=\binom{k}{r-1}
\).
Then
\[
  \Delta q(k)
  =4\left(\binom{k+1}{4}-\binom{k}{4}\right)
  =4\binom{k}{3}.
\]
Applying \(\Delta\) repeatedly gives
\[
  \Delta^2 q(k)=4\binom{k}{2},
\]
\[
  \Delta^3 q(k)=4\binom{k}{1}=4k,
\]
\[
  \Delta^4 q(k)=4\binom{k}{0}=4,
\]
\[
  \Delta^5 q(k)=0.
\]
\end{proof}

\begin{corollary}[Order--5 linear recurrence]
\label{cor:numbers-order5-recurrence}
The sequence \(q(k)\) satisfies
\[
  q(k+5)-5q(k+4)+10q(k+3)-10q(k+2)+5q(k+1)-q(k)=0.
\]
\end{corollary}

\begin{proof}
The left-hand side is exactly \(\Delta^5 q(k)\), which vanishes by
\cref{thm:numbers-difference-tower}.
\end{proof}

\begin{remark}[Rigidity viewpoint]
\label{rem:numbers-rigidity-delta4}
The identity \(\Delta^4 q\equiv4\) characterizes \(q\) among sequences with
zero initial values at depths \(0,1,2,3\):
the quartic law is the unique depth-polynomial with constant fourth
difference equal to \(4\).
\end{remark}

% ============================================================
\section{Monotonicity, convexity, and order transfer}
\label{sec:numbers-order-transfer}
% ============================================================

\begin{proposition}[Strict growth]
\label{prop:numbers-strict-growth}
For all \(k\ge4\),
\[
  q(k+1)-q(k)=4\binom{k}{3}>0.
\]
Hence \(q\) is strictly increasing on \(\{4,5,6,\dots\}\).
\end{proposition}

\begin{proof}
The difference formula is \cref{thm:numbers-difference-tower}.
For \(k\ge4\), \(\binom{k}{3}>0\), so the increment is positive.
\end{proof}

\begin{proposition}[Discrete convexity]
\label{prop:numbers-discrete-convexity}
For all \(k\ge2\),
\[
  \Delta^2 q(k)=4\binom{k}{2}>0.
\]
Thus the increments \(q(k+1)-q(k)\) are strictly increasing.
\end{proposition}

\begin{proof}
By \cref{thm:numbers-difference-tower},
\(
  \Delta^2 q(k)=4\binom{k}{2}
\).
For \(k\ge2\), \(\binom{k}{2}>0\).
Therefore \(\Delta q\) is strictly increasing.
\end{proof}

\begin{theorem}[Depth order transfers to mass order]
\label{thm:numbers-depth-order-transfer}
For the closed-system mass coordinate of \cref{thm:numbers-mass-calibration}, let
\(\sigma,\tau\) be nontrivial primitive sectors with
\(k_\sigma,k_\tau\ge4\).
If
\[
  k_\sigma<k_\tau,
\]
then
\[
  m_\sigma<m_\tau.
\]
\end{theorem}

\begin{proof}
By \cref{cor:numbers-conditional-mass-law},
\(
  m_x=\mu q(k_x)
\)
for \(x\in\{\sigma,\tau\}\), with \(\mu>0\).
By \cref{prop:numbers-strict-growth},
\(q(k_\sigma)<q(k_\tau)\).
Multiplication by positive \(\mu\) preserves strict inequality.
\end{proof}

\begin{corollary}[Injectivity on realized depth values]
\label{cor:numbers-depth-mass-injective}
For the closed-system mass coordinate, the map
\(
  k\mapsto m(k):=4\mu\binom{k}{4}
\)
is injective on \(\{4,5,6,\dots\}\).
\end{corollary}

\begin{proof}
If \(m(k_1)=m(k_2)\), strict monotonicity from
\cref{prop:numbers-strict-growth} implies \(k_1=k_2\).
\end{proof}

% ============================================================
\section{Inversion bounds for the mass-coordinate law}
\label{sec:numbers-inversion}
% ============================================================

\begin{proposition}[Elementary quartic bounds for \(\binom{k}{4}\)]
\label{prop:numbers-elementary-bounds}
For every integer \(k\ge4\),
\[
  (k-3)^4
  \le
  k(k-1)(k-2)(k-3)
  \le
  k^4.
\]
Equivalently,
\[
  \frac{(k-3)^4}{24}
  \le
  \binom{k}{4}
  \le
  \frac{k^4}{24}.
\]
\end{proposition}

\begin{proof}
For \(k\ge4\), each factor in
\(k(k-1)(k-2)(k-3)\)
lies between \(k-3\) and \(k\), so the product lies between the
corresponding fourth powers.
Division by \(24\) gives the binomial form.
\end{proof}

\begin{theorem}[Depth window from one mass value]
\label{thm:numbers-depth-window}
For the closed-system mass coordinate of \cref{thm:numbers-mass-calibration}, let
\(m>0\) satisfy
\[
  m=4\mu\binom{k}{4}
\]
for some integer \(k\ge4\), and define
\[
  x:=\left(\frac{6m}{\mu}\right)^{1/4}.
\]
Then
\[
  k-3\le x\le k,
\]
hence
\[
  x\le k\le x+3.
\]
\end{theorem}

\begin{proof}
From the mass law,
\[
  \frac{6m}{\mu}=k(k-1)(k-2)(k-3).
\]
Apply \cref{prop:numbers-elementary-bounds}:
\[
  (k-3)^4\le \frac{6m}{\mu}\le k^4.
\]
Taking fourth roots gives
\(k-3\le x\le k\), and rearrangement gives
\(x\le k\le x+3\).
\end{proof}

\begin{corollary}[Finite-candidate inversion algorithm]
\label{cor:numbers-finite-candidate-inversion}
For the closed-system mass coordinate, a mass-coordinate value
determines depth by checking at most four consecutive integers.
\end{corollary}

\begin{proof}
By \cref{thm:numbers-depth-window}, any admissible integer depth must lie in
\([x,x+3]\), which contains at most four integers.
By \cref{prop:numbers-strict-growth}, at most one of these candidates can
match the observed mass.
\end{proof}

\begin{proposition}[Asymptotic depth extraction]
\label{prop:numbers-asymptotic-depth}
For the closed-system mass coordinate, for large depth values one has
\[
  k
  =
  \left(\frac{6m}{\mu}\right)^{1/4}
  +O(1).
\]
\end{proposition}

\begin{proof}
This is immediate from \cref{thm:numbers-depth-window}.
The difference between \(k\) and the quartic root is always bounded by
\(3\).
\end{proof}

% ============================================================
\section{Coupling to the forced quadratic normalization}
\label{sec:numbers-quadratic-coupling}
% ============================================================

\begin{definition}[Quadratic and quartic carriers at depth \(k\)]
\label{def:numbers-M2-M4}
Define
\[
  M_2(k):=\binom{k}{2},
\]
\[
  M_4(k):=4\mu\binom{k}{4}.
\]
Here \(M_2\) is the structural quadratic count
(\cref{thm:binomk2}), and \(M_4\) is the quartic mass-coordinate law.
\end{definition}

\begin{proposition}[Exact ratio formula \(M_4/M_2^2\)]
\label{prop:numbers-ratio-M4-M2}
For every \(k\ge4\),
\[
  \frac{M_4(k)}{M_2(k)^2}
  =
  \frac{2\mu}{3}
  \cdot
  \frac{(k-2)(k-3)}{k(k-1)}.
\]
\end{proposition}

\begin{proof}
Using \cref{def:numbers-M2-M4},
\[
  \frac{M_4(k)}{M_2(k)^2}
  =
  \frac{4\mu\binom{k}{4}}{\binom{k}{2}^2}
  =
  4\mu
  \frac{\frac{k(k-1)(k-2)(k-3)}{24}}
       {\frac{k^2(k-1)^2}{4}}.
\]
Simplifying yields
\[
  \frac{2\mu}{3}\frac{(k-2)(k-3)}{k(k-1)}.
\]
\end{proof}

\begin{corollary}[Asymptotic quadratic coupling]
\label{cor:numbers-asymptotic-quadratic-coupling}
As \(k\to\infty\),
\[
  M_4(k)
  \sim
  \frac{2\mu}{3}\,M_2(k)^2.
\]
\end{corollary}

\begin{proof}
In \cref{prop:numbers-ratio-M4-M2}, the rational factor
\(
  \frac{(k-2)(k-3)}{k(k-1)}
\)
tends to \(1\).
Therefore
\(
  M_4/M_2^2\to 2\mu/3
\).
\end{proof}

\begin{remark}[Interpretive consequence]
\label{rem:numbers-coupling-interpretation}
Under the same depth variable, the quartic law is asymptotically quadratic
in the quadratic carrier count.
Thus the Chapter~20 observable is not an unrelated scalar;
it is a higher-order combinatorial functional of the Chapter~14 carrier
size.
\end{remark}

% ============================================================
\section{Filtration-time scaling under forced normalization}
\label{sec:numbers-time-scaling}
% ============================================================

\begin{definition}[Time-lifted quartic observable]
\label{def:numbers-time-lifted}
Let \(k(T)\) denote activated depth at filtration time \(T\), as in
\cref{thm:forced-normalization}.
Define
\[
  M_4(T):=4\mu\binom{k(T)}{4}.
\]
\end{definition}

\begin{theorem}[Quadratic-time scaling of the quartic observable]
\label{thm:numbers-time-quadratic}
Under the canonical normalization of
\cref{thm:forced-normalization,lem:parabolic}, one has
\[
  M_4(T)\asymp T^2.
\]
\end{theorem}

\begin{proof}
By \cref{lem:parabolic},
\(
  k(T)\asymp\sqrt T
\).
By \cref{prop:numbers-qk-polynomial},
\(
  \binom{k}{4}\asymp k^4
\)
for large \(k\).
Therefore
\[
  M_4(T)
  =4\mu\binom{k(T)}{4}
  \asymp
  k(T)^4
  \asymp
  (\sqrt T)^4
  =T^2.
\]
\end{proof}

\begin{corollary}[Separation of visible scales]
\label{cor:numbers-visible-scale-separation}
With
\(
  M_2(T)\asymp T
\)
from \cref{thm:forced-normalization}, the quartic observable satisfies
\[
  M_4(T)\asymp M_2(T)^2.
\]
\end{corollary}

\begin{proof}
Combine \cref{thm:forced-normalization} with
\cref{thm:numbers-time-quadratic}.
\end{proof}

% ============================================================
\section{General anchor formalism and ratio geometry}
\label{sec:numbers-anchor-geometry}
% ============================================================

\begin{definition}[Anchor package]
\label{def:numbers-anchor-package}
An anchor package is a tuple
\[
  \mathcal A=(k_1,\dots,k_r),
\]
with integers \(k_i\ge4\), together with the mass-coordinate
bridge theorem \cref{thm:numbers-mass-calibration}.
\end{definition}

\begin{theorem}[General ratio law]
\label{thm:numbers-general-ratio-law}
Under \cref{def:numbers-anchor-package}, for any indices \(i,j\),
\[
  \frac{m_i}{m_j}
  =
  \frac{\binom{k_i}{4}}{\binom{k_j}{4}}.
\]
\end{theorem}

\begin{proof}
By \cref{cor:numbers-conditional-mass-law},
\(
  m_i=4\mu\binom{k_i}{4}
\)
and similarly for \(m_j\).
Divide and cancel the common factor \(4\mu\).
\end{proof}

\begin{proposition}[Multiplicative transitivity of ratio data]
\label{prop:numbers-ratio-transitivity}
For any \(i,j,\ell\),
\[
  \frac{m_i}{m_j}\cdot\frac{m_j}{m_\ell}=\frac{m_i}{m_\ell}.
\]
On the depth side this reads
\[
  \frac{\binom{k_i}{4}}{\binom{k_j}{4}}
  \cdot
  \frac{\binom{k_j}{4}}{\binom{k_\ell}{4}}
  =
  \frac{\binom{k_i}{4}}{\binom{k_\ell}{4}}.
\]
\end{proposition}

\begin{proof}
This is immediate in the multiplicative group \(\mathbb R_{>0}\),
or directly from \cref{thm:numbers-general-ratio-law}.
\end{proof}

\begin{corollary}[Recovery of the \((6,18,35)\) package]
\label{cor:numbers-anchor-61835}
For the explicit depth-anchor package of \cref{def:numbers-depth-anchors},
\[
  \frac{m_\mu}{m_e}=\frac{\binom{18}{4}}{\binom{6}{4}}=204,
\]
\[
  \frac{m_\tau}{m_e}=\frac{\binom{35}{4}}{\binom{6}{4}}=\frac{10472}{3}.
\]
\end{corollary}

\begin{proof}
Apply \cref{thm:numbers-general-ratio-law} to
\((k_e,k_\mu,k_\tau)=(6,18,35)\), then evaluate binomial coefficients.
\end{proof}

The explicit \((6,18,35)\) anchor package can also be recovered from
ratio data once third-generation onset depth is fixed by the cubic
selector route.
Dependency basis:
\cref{thm:numbers-general-ratio-law,prop:numbers-strict-growth,cor:numbers-k-tau-35}.

\begin{theorem}[Ratio/cubic closure of the \texorpdfstring{$(6,18,35)$}{(6,18,35)} anchors]
\label{thm:numbers-anchor-closure-61835}
For the closed-system mass coordinate of \cref{thm:numbers-mass-calibration}, let
\(e,\mu,\tau\) be reference sectors with depths
\(k_e,k_\mu,k_\tau\ge4\).
Suppose the following calibration data hold:
\begin{enumerate}[label=\textup{(\roman*)}]
\item
      \(k_\tau=35\) (for instance by \cref{cor:numbers-k-tau-35});
\item
      measured mass ratios satisfy
      \[
        \frac{m_\mu}{m_e}=204,
        \qquad
        \frac{m_\tau}{m_e}=\frac{10472}{3}.
      \]
\end{enumerate}
Then
\[
  k_e=6,
  \qquad
  k_\mu=18,
  \qquad
  k_\tau=35.
\]
In particular, the explicit depth-anchor package of
\cref{def:numbers-depth-anchors} is recovered.
\end{theorem}

\begin{proof}
From \cref{thm:numbers-general-ratio-law} and \(k_\tau=35\),
\[
  \frac{\binom{35}{4}}{\binom{k_e}{4}}
  =
  \frac{m_\tau}{m_e}
  =
  \frac{10472}{3}.
\]
Hence
\[
  \binom{k_e}{4}
  =
  \frac{\binom{35}{4}}{10472/3}
  =
  15
  =
  \binom{6}{4}.
\]
By \cref{prop:numbers-strict-growth},
\(k\mapsto q(k)=4\binom{k}{4}\) is strictly increasing on
\(\{4,5,6,\dots\}\), so \(k\mapsto\binom{k}{4}\) is strictly increasing
there; therefore \(k_e=6\).

Applying \cref{thm:numbers-general-ratio-law} to \((\mu,e)\),
\[
  \frac{\binom{k_\mu}{4}}{\binom{6}{4}}
  =
  \frac{m_\mu}{m_e}
  =
  204,
\]
so
\[
  \binom{k_\mu}{4}=204\cdot\binom{6}{4}=204\cdot15=3060=\binom{18}{4}.
\]
Strict monotonicity again yields \(k_\mu=18\), and \(k_\tau=35\) is the
given cubic-onset depth.
Thus \((k_e,k_\mu,k_\tau)=(6,18,35)\), i.e. the explicit
depth-anchor package of \cref{def:numbers-depth-anchors} is recovered.
\end{proof}

\begin{corollary}[Discharge of explicit depth-anchor package]
\label{cor:numbers-anchor-closure-discharges}
Whenever \cref{thm:numbers-anchor-closure-61835} applies, consequences attached to the explicit depth-anchor package
of \cref{def:numbers-depth-anchors} can be used without stipulating that
package independently.
\end{corollary}

\begin{proof}
Immediate from \cref{thm:numbers-anchor-closure-61835}.
\end{proof}

% ============================================================
\section{Local ratio sensitivity and depth identifiability}
\label{sec:numbers-ratio-sensitivity}
% ============================================================

\begin{proposition}[One-step ratio formula]
\label{prop:numbers-one-step-ratio}
For integer \(k\ge4\),
\[
  \frac{m(k+1)}{m(k)}
  =
  \frac{\binom{k+1}{4}}{\binom{k}{4}}
  =
  \frac{k+1}{k-3},
\]
where \(m(k):=4\mu\binom{k}{4}\).
\end{proposition}

\begin{proof}
Compute directly:
\[
  \frac{\binom{k+1}{4}}{\binom{k}{4}}
  =
  \frac{(k+1)k(k-1)(k-2)}{k(k-1)(k-2)(k-3)}
  =
  \frac{k+1}{k-3}.
\]
\end{proof}

\begin{corollary}[Closed-form depth from successive ratio]
\label{cor:numbers-depth-from-successive-ratio}
If a measured consecutive ratio is
\[
  r:=\frac{m(k+1)}{m(k)}>1,
\]
then
\[
  k=\frac{3r+1}{r-1}.
\]
\end{corollary}

\begin{proof}
By \cref{prop:numbers-one-step-ratio},
\(
  r(k-3)=k+1
\).
Rearrange:
\(
  k(r-1)=3r+1
\), hence the formula.
\end{proof}

\begin{proposition}[Relative sensitivity to a \(\pm1\) depth shift]
\label{prop:numbers-pm1-sensitivity}
For \(k\ge5\),
\[
  \frac{m(k+1)-m(k)}{m(k)}=\frac{4}{k-3},
\]
\[
  \frac{m(k)-m(k-1)}{m(k)}=\frac{4}{k}.
\]
Thus the relative one-step sensitivity decays like \(O(1/k)\).
\end{proposition}

\begin{proof}
From \cref{prop:numbers-one-step-ratio},
\[
  \frac{m(k+1)-m(k)}{m(k)}
  =
  \frac{m(k+1)}{m(k)}-1
  =
  \frac{k+1}{k-3}-1
  =
  \frac{4}{k-3}.
\]
Similarly,
\[
  \frac{m(k-1)}{m(k)}
  =
  \frac{\binom{k-1}{4}}{\binom{k}{4}}
  =
  \frac{k-4}{k},
\]
so
\[
  \frac{m(k)-m(k-1)}{m(k)}
  =
  1-\frac{k-4}{k}
  =
  \frac{4}{k}.
\]
\end{proof}

\begin{remark}[Practical calibration implication]
\label{rem:numbers-calibration-implication}
At larger anchor depths, fixed absolute uncertainty in depth produces
smaller relative distortion in mass ratios.
This is a direct consequence of
\cref{prop:numbers-pm1-sensitivity}.
\end{remark}

% ============================================================
\section{Cubic selector in symmetric-polynomial normal form}
\label{sec:numbers-cubic-normal-form}
% ============================================================

\begin{definition}[Power sums and elementary sums]
\label{def:numbers-power-elementary}
For charges \(q_1,\dots,q_k\), define
\[
  p_r:=\sum_{i=1}^k q_i^r
  \qquad (r\ge1),
\]
and denote by \(e_r\) the elementary symmetric polynomials.
In this notation,
\[
  p_1=e_1,
  \qquad
  e_3=\sum_{1\le i<j<\ell\le k}q_iq_jq_\ell.
\]
\end{definition}

\begin{proposition}[Neutrality identity]
\label{prop:numbers-neutrality-identity}
If \(p_1=0\), then
\[
  p_3=3e_3.
\]
In particular,
\[
  \sum_{i=1}^k q_i^3
  =
  3\sum_{1\le i<j<\ell\le k}q_iq_jq_\ell.
\]
\end{proposition}

\begin{proof}
Newton's identity gives
\[
  p_3=e_1p_2-e_2p_1+3e_3.
\]
Since \(p_1=e_1=0\), the first two terms vanish, yielding
\(p_3=3e_3\).
\end{proof}

\begin{theorem}[Equivalent cubic selector forms]
\label{thm:numbers-cubic-equivalent-forms}
Assume charge neutrality
\(
  \sum_i q_i=0
\).
Then the following are equivalent:
\begin{enumerate}[label=\textup{(\roman*)}]
\item
      \(\sum_{i<j<\ell}q_iq_jq_\ell=0\);
\item
      \(\sum_i q_i^3=0\);
\item
      \(\Delta_3(k)=0\).
\end{enumerate}
\end{theorem}

\begin{proof}
By \cref{prop:numbers-neutrality-identity},
\(
  \sum_i q_i^3=3\sum_{i<j<\ell}q_iq_jq_\ell
\).
Thus (i) and (ii) are equivalent.
Condition (iii) is exactly the notation of (ii) by
\cref{def:numbers-cubic-defect}.
\end{proof}

\begin{proposition}[General minimal-depth criterion]
\label{prop:numbers-cubic-minimal-general}
Suppose admissibility requires
\(\Delta_3(k)=0\), and there exists \(k_\ast\ge4\) such that
\[
  \Delta_3(k)\neq0
  \quad\text{for all }4\le k<k_\ast,
\]
\[
  \Delta_3(k_\ast)=0.
\]
Then the minimal admissible depth is \(k_\ast\).
\end{proposition}

\begin{proof}
The first line excludes every depth smaller than \(k_\ast\).
The second line certifies admissibility at \(k_\ast\).
Therefore no admissible depth exists below \(k_\ast\), and \(k_\ast\)
is the least admissible depth.
\end{proof}

\begin{corollary}[Specialization to \(k_\tau=35\)]
\label{cor:numbers-k35-specialization}
If
\[
  \Delta_3(k)\neq0\ \text{for all }4\le k\le34,
\]
\[
  \Delta_3(35)=0,
\]
then the first cubic-admissible depth is \(35\).
\end{corollary}

\begin{proof}
Apply \cref{prop:numbers-cubic-minimal-general} with
\(k_\ast=35\).
\end{proof}

% ============================================================
\section{Integrated conditional closure theorem}
\label{sec:numbers-integrated-closure}
% ============================================================

\begin{theorem}[Integrated conditional numerical closure]
\label{thm:numbers-integrated-closure}
Fix \(\mu>0\) and a primitive sector assignment
\(\sigma\mapsto k_\sigma\in\mathbb N\).
Suppose the following data are fixed:
\begin{enumerate}[label=\textup{(\roman*)}]
\item
      structural depth and carrier data from
      \cref{def:numbers-dependency-interface};
\item
      optional anchor/depth calibration route: either explicit anchors
      \cref{def:numbers-depth-anchors} or the ratio/cubic route
      \cref{thm:numbers-anchor-closure-61835};
\item
      optional cubic selector package
      \cref{def:numbers-phase-admissible,cor:numbers-cubic-criterion}.
\end{enumerate}
Then:
\begin{enumerate}[label=\textup{(\alph*)}]
\item
      \(m_\sigma=4\mu\binom{k_\sigma}{4}\) for every primitive sector;
\item
      mass ordering follows depth ordering for all depths \(\ge4\);
\item
      all anchor ratios are given by binomial quotients
      \(\binom{k_i}{4}/\binom{k_j}{4}\);
\item
      if the cubic defect first vanishes at depth \(k_\ast\), then
      \(k_\ast\) is the first cubic-admissible generation depth;
\item
      if either the explicit depth-anchor package of
      \cref{def:numbers-depth-anchors} is stipulated or
      \cref{thm:numbers-anchor-closure-61835} applies, then the concrete
      \((6,18,35)\) depth-anchor package is fixed.
\end{enumerate}
\end{theorem}

\begin{proof}
Part (a) is \cref{cor:numbers-conditional-mass-law}, using
\cref{thm:numbers-mass-calibration}.
Part (b) is \cref{thm:numbers-depth-order-transfer}.
Part (c) is \cref{thm:numbers-general-ratio-law}.
Part (d) is \cref{prop:numbers-cubic-minimal-general}.
Part (e) is immediate from
\cref{def:numbers-depth-anchors,thm:numbers-anchor-closure-61835}.
All components are conditional exactly on the stated structural and
calibration data.
\end{proof}

\begin{corollary}[Operational checklist]
\label{cor:numbers-operational-checklist}
To apply Chapter~20 numerically, it is sufficient to
specify:
\begin{enumerate}[label=\textup{(\roman*)}]
\item
      the depth assignment \(k_\sigma\);
\item
      the normalization constant \(\mu\);
\item
      optional anchor/depth calibration route: either explicit reference
      depths, or the ratio/cubic route of
      \cref{thm:numbers-anchor-closure-61835};
\item
      optional cubic selector data \(\Delta_3(k)\)
      (used in particular for the ratio/cubic route).
\end{enumerate}
After these data are fixed, every Chapter~20 output is an explicit algebraic
consequence.
\end{corollary}

\begin{proof}
This is a direct restatement of
\cref{thm:numbers-integrated-closure} and the formulas proved in
\cref{sec:numbers-discrete-calculus,sec:numbers-anchor-geometry,sec:numbers-cubic-normal-form}.
\end{proof}

\begin{remark}[Why this chapter is conditionally complete]
\label{rem:numbers-conditionally-complete}
From the structural side, multilinear combinatorics on
\(\Lambda^2(\mathbb Z^k)\) fixes the exact quartic backbone
\(E_4(\omega_k)=4\binom{k}{4}\), while the broader admissible-depth
rigidity package fixes quartic growth and quartic leading-order
behavior at the stated theorem strength.
From the quantitative side, the remaining degrees of freedom are
named explicitly as conditional boundary data, and the formulas
attached to each chosen data package display those dependencies
directly.
This is the precise sense in which Chapter~20 is complete as a
\emph{conditional} numerical closure theorem.
\end{remark}

% ============================================================
\section{Conclusion}
\label{sec:numbers-conclusion}
% ============================================================

The chapter now yields a full dependency-explicit numerical package with
three layers:
\begin{enumerate}[label=\textup{(\roman*)}]
\item The quartic combinatorics are exact and internally complete:
      \(E_4(\omega_k)=4\binom{k}{4}\)
      (\cref{prop:numbers-quartic-count,cor:numbers-energy-by-counting}).
\item The induced depth polynomial has explicit discrete calculus,
      recurrence, monotonicity, convexity, and inversion theory
      (\cref{thm:numbers-difference-tower,cor:numbers-order5-recurrence,thm:numbers-depth-window}).
\item The quartic mass-coordinate law is identified by the intrinsic bridge theorem:
      \(m_\sigma=4\mu\binom{k_\sigma}{4}\)
      (\cref{thm:numbers-mass-calibration,cor:numbers-conditional-mass-law}).
\item Within a fixed anchor package, anchor ratios are given exactly
      by binomial quotients
      (\cref{thm:numbers-general-ratio-law,cor:numbers-anchor-61835}).
\item The explicit \((6,18,35)\) anchor package can be recovered via
      ratio/cubic closure and therefore need not be stipulated separately
      when that route is available
      (\cref{thm:numbers-anchor-closure-61835,cor:numbers-anchor-closure-discharges}).
\item The cubic selector admits symmetric-polynomial normal forms and a
      general minimal-depth theorem
      (\cref{thm:numbers-cubic-equivalent-forms,prop:numbers-cubic-minimal-general}).
\item The full package is assembled as one conditional closure theorem
      (\cref{thm:numbers-integrated-closure,cor:numbers-operational-checklist}).
\end{enumerate}

Accordingly, Chapter~20 now aligns with the proof architecture used in the
longer structural chapters:
it provides an exact quartic combinatorial core,
an explicit dependency ledger,
and a conditional chain from stated boundary data to the corresponding
numerical outputs.

What remains beyond this chapter is precisely the status of
in-shell splitting and final normalization: whether an additional
intrinsic shell invariant is needed, whether its coefficient is
internally fixed or irreducible, and how absolute mass scale is anchored.
That extension is isolated in \cref{ch:numbers-shell-closure}.

\begin{chapterclosing}
\closingitem{Established}{The quartic combinatorial core, mass-depth law, anchor-ratio algebra, and cubic-selector package are assembled as a conditional numerical closure theorem.}
\closingitem{Carried forward}{The remaining numerical question is the status of shell-level splitting, correction structure, and normalization.}
\end{chapterclosing}
